%% Dans ce fichier, les �ventuels accents sont cod�s
%% en ISO-8859-15 (aka latin9)
%% La ligne suivante donne trois exemples 
%% e accent aigu : �, e accent grave : �, e-dans-l'o : �
%%%%%ENTREES%%%%
\begin{tabarg}
\multicolumn{3}{l}{\titretab {Entr\'ees}}\\
\hline
\$Arg1 & Coeff., \texttt{float} & blabla (m)\\
\$Arg2 & Coeff., \texttt{float} & blabla (m)\\
\$Arg3 & Coeff., \texttt{float} & blabla (m)\\
\$Arg4 & Vecteur, \texttt{float} & blabla (m/s)\\
\$Arg5 & Vecteur, \texttt{float} & blabla (m/s)\\
\$Arg6 & Vecteur, \texttt{float} & blabla (m/s)\\
\$Arg7 & Vecteur, \texttt{float} & blabla (m/s)\\
\$Arg8 & Vecteur, \texttt{float} & blabla (m/s)\\
\$Arg9 & Vecteur, \texttt{float} & blabla (m/s)\\
\$Arg10 & Vecteur, \texttt{float} & blabla (radian/s)\\
\$Arg11 & Vecteur, \texttt{float} & blabla (radian/s)\\
\$Arg12 & Vecteur, \texttt{float} & blabla (radian/s)\\
\$Arg13 & Vecteur, \texttt{float} & blabla (radian)\\
\$Arg14 & Vecteur, \texttt{float} & blabla (radian)\\
\$Arg15 & Vecteur, \texttt{float} & blabla (radian)\\
\end{tabarg}

\medskip
%%%%%SORTIES%%%%
\begin{tabarg}
\multicolumn{3}{l}{\titretab {Sorties}}\\
\hline
\$Arg16 & Vecteur, \texttt{float} & blabla (\$Arg4.unit)\\
\$Arg17 & Vecteur, \texttt{float} & blabla (\$Arg4.unit)\\
\$Arg18 & Vecteur, \texttt{float} & blabla (\$Arg4.unit)\\
\end{tabarg}
